\section{ベイズ推測の基礎}
\label{D1_lesson25}
\subsection{授業内容}
\subsubsection{科目の中でのこのコマの位置づけ}
\begin{tabular}[t]{cccccc}
    記述統計[4] & 確率[12] & モデル[16] & 推測・推定[7] & 意思決定[6] & 実践[4]
\end{tabular}
\subsubsection{概要}
母数を推定する第三の方法として，ベイズ法による推定を考える。ベイズの定理の歴史は古く，逆確率の法則とも呼ばれるものである。ベイズの定理を理解するにあたっては，条件つき確率について理解する必要があり，これらを通じてベイズの定理を導く。ベイズ法の利点は，尤度関数を確率分布関数に変えてくれることであり，推定値の確からしさを確率で表現することができる点である。ここでは病気の罹患率の例を用いて説明し，確率を信念の更新として理解することもできることを学ぶ。
\subsubsection{コマ主題細目(箇条書き)}
\begin{description}
\item[二次元の確率変数]これまで扱ってきた確率分布は単一のものであったが，二変数による(二次元的な)確率分布を導入する。複数の事象が同時に生じることを表す同時確率，同時確率のうち一方の変数だけに周辺化して考える周辺確率，一方の情報が分かっている上での確率である条件つき確率，の3つの考え方が導かれる。また，両者が独立であるとはどういうことかについて確率的な定義を与える。簡単な数値例を用いて解説する。
\begin{flushright}
    $\to$ \citet{Maeda2017}のP.93--109.\citet{豊田201506}のP.3--11.
\end{flushright}

\item[ベイズの定理]同時確率，周辺確率の定義などを用いて，条件付き確率の式を変形していくことができ，最終的に右辺と左辺で条件が逆になる場合の確率が計算できることがわかる。これがベイズの定理と呼ばれるものであり，逆確率の定理，法則と呼ばれることもある。ベイズの定理は，事前確率に尤度と周辺尤度の比をかけたものとして理解することもできるし，尤度関数を確率として捉えるために変形させる項として事前確率と周辺尤度の比をかけたもの，と考えることもできる。
\begin{flushright}
    $\to$ \citet{Maeda2017}のP.93--109.尤度を確率分布に変えることの分かりやすい説明としては，\citet{Lambert201805}のTable4.1.Xが良い。
\end{flushright}

\item[推論と確率]ベイズの定理は確率のアップデートとして考えることができる。我々が推論するというのは，確率をアップデートすることだと考えることもできるのであれば。確率という数字が従うべきルールさえ守られていれば，確率は人間の自然な推論モデルとして考えることができるだろう。ここでは病気の罹患率の例を使って，その計算手順と，事前分布と事後分布の関係について理解する。
\begin{flushright}
    $\to$ \citet{Maeda2017}のP.15--23.あるいは\citet{豊田201506}のP.8--12.
\end{flushright}

\end{description}
\subsubsection{キーワード}
\begin{itemize}
\item 同時確率
\item 周辺確率
\item 条件付き確率
\item ベイズの定理
\item 事前確率，事後確率
\end{itemize}
\subsection{授業情報}
\paragraph{コマの展開方法}
講義
\paragraph{標準シラバスにおける位置づけ}
科目番号5；心理学統計法；(1)心理学で用いられる統計手法；J より高度な記述や推定を目指して
\subsubsection{予習・復習課題}
\paragraph{予習}
ここからは，より進んだ統計的内容に入る。ベイズ統計学を心理学領域で活用する歴史はまだ浅いが，現代社会では実践的に多用されている技術であり，心理学にとっても今後ますます重要性が高まることは間違いない。例えば\citet{野島201903}は公認心理師向けのテキストであるが，その最終章をベイズにあてている。
内容的には，確率についてまた違った視点からの考え方が導入される。とはいえ，確率が守るべきルールについては変わりがない。改めて第\ref{D1_lesson06}講の確率について復習しておくとよい。
「ベイズ統計学」をキーワードとしたいくつかの入門書があるので，目を通しておくと良い。あるいは「モンティ・ホール問題」で検索すると，興味を覚える人もいるかもしれない。
\paragraph{復習}
さまざまな確率の書き方から，ベイズの定理の導出まで，筋道をもう一度自らフォローアップして導出できるかどうか確認しておくこと。罹患率のような具体例については，読み物として\citet{Gerd200309}があるので目を通すと良い。
確率についての説明の仕方については，\citet{平岡200910}のP.1--69の説明も異なる切り口からで面白いし，\citet{石田201801}のP.1--12もわかりやすい。ベイズ統計学についてのごく簡単な入門書としては，ラノベ風の\citet{石田石田201902}なども良いだろう。